\section{The fixed-tail entrance at $u=h+1$}
Fix a nonempty sorted word $V=(t_1,\ldots,t_r)$ with $t_i\ge3$.
Put $W=\sum t_i$, $h=W-r$, $H=h+1$, and
$F_{u,\alpha}=F_{1^\alpha2^uV}$, $\alpha\ge1$.
Let $D$ be its actual positive slope from rank two, and let $\ell$
be the FIRST global $D$ mode. This result strengthens the previously
proved entrance $h+4$; the new entrance remains an explicit hypothesis.

\begin{theorem}[Sharpened single-entrance finite principle]
Suppose at the SINGLE binary count $u=H=h+1$, for EVERY
$1\le\alpha\le H+h+1=2h+2$, the actual slopes are LC at all centres
\[
                         H+r+1=W+2\le n<\ell.           \tag{1}
\]
Then for ALL $u\ge H$, ALL $\alpha\ge1$ and ALL current ranks
$N\ge W+1$, every actual C comparison holds.
If the reference additionally satisfies C at its finite earlier
ranks $N\le W$, complete actual C follows for all $u\ge H$.
If, in addition, complete actual C is checked for
$0\le u<H$, $1\le\alpha\le u+h+1$, it follows for ALL $u\ge0$.
The reference and actual finite tests use coefficients only through
$W$ and $2u+W+1$, respectively. No assertion that every entrance
passes, or that general original unimodality follows, is included.
\end{theorem}

\subsection*{The alternating boundary pays LC from $h+1$}
Use the exact reference $F^\star$, slopes $r_n$, and positive
Cassini amplitude $A_V$ of the earlier fixed-tail principle.
Its signed homogeneous minor, INCLUDING its sign, is
\[
       r_n^2-r_{n-1}r_{n+1}=(-1)^{n-W}A_V\quad(n\ge W+1).
                                                               \tag{2}
Indeed $A_V=|B(q)B(-\varphi)|$ and
$\operatorname{sgn}B(-\varphi)=(-1)^W$ in the Binet determinant.
The old entire-word budget remains $r_{W+2}\ge A_V$ and
$r_{W+4}\ge2A_V+1$.

The first actual defect is at $N_0=u+r+2$ and equals
$\delta=2$ for $\alpha=1$, and one otherwise.
At $n=N_0-1=u+r+1\ge W+2$ the actual boundary minor is exactly
\[
                       (-1)^{n-W}A_V+\delta r_{n-1}.    \tag{3}
When $n-W$ is even it is positive. When odd, $n\ge W+3$, so
$r_{n-1}\ge r_{W+2}\ge A_V$ and (3) is nonnegative.
Only nonnegativity is required; a zero minor is allowed.

Now use the SAME whole polynomial insertion identity
$F_{u+1,\alpha}=S_2F_{u,\alpha}+B_V$.
Its slope forcing has degree $W+1$. The earlier proof's two-term
LC propagation is valid above this degree, uses old LC only through
the FIRST old mode, and uses only a nonnegative new boundary minor.
The new centres all lie at or beyond $W+2$. The unique possible
pre-mode flat is at lastcapacity$+1\le W+1$, outside this region.
Thus the same induction proves LC at all $u+r+1\le n\le\ell$
for EVERY $u\ge h+1$. Padding clamps and mode/read bounds are the
unchanged exact arguments of the earlier finite principle.

\subsection*{The replacement for the unnecessarily early D seeds}
The new general reference theorem proves ALL C shifts from $W+1$.
Exact prefix agreement pays actual ranks $W+1\le N<N_0$;
at $N_0$ the first defect only improves C.
For a later strictly increasing actual current slope, put
$\lambda=D_N/D_{N-1}>1$ and
$\lambda_0=(r_{N_0}-\delta)/r_{N_0-1}$.
Propagated LC gives $\lambda\le\lambda_0$ and all required later
$D_i\ge\lambda D_{i-1}$ starting at $i=N_0$.

Write $(b,c,d)=(r_{N_0-2},r_{N_0-1},r_{N_0})$, so $d=b+c$.
Let $k_0=N_0-3$. Reference C at $N_0$, together with $\lambda\le
\lambda_0$, pays $G_k\ge\lambda G_{k-1}$ at EVERY $k\le k_0$.
The ordinary coefficients at the next two reads are exactly
\[
                F_{k_0}=c,\quad F_{k_0+1}=d,
                         \quad F_{k_0+2}=c+d.           \tag{4}
\]
All three ranks precede the first actual defect. The homogeneous
identity $F^star_j=r_{j+2}$ applies there: the smallest read
$k_0=N_0-3\ge W$.
The first ordinary comparison $d\ge\lambda c$ is immediate.
For the second,
\[
 c(c+d)-d(d-\delta)=c^2-bd+\delta d
                          =(-1)^{N_0-1-W}A_V+\delta d\ge0. \tag{5}
\]
If the Cassini sign is positive this is immediate. If negative,
$N_0-1\ge W+3$, hence $d\ge r_{W+4}\ge2A_V+1$.
Thus $F_{k_0+2}\ge\lambda F_{k_0+1}$ as well.
Using $G_{j}=G_{j-1}+F_j$, these two ordinary comparisons extend
the cumulative comparison twice, to $k_0+1$ and $k_0+2$.

For every subsequent required rank use
$F_i=F_{i-1}+D_i$ and $G_i=G_{i-1}+F_i$.
The $F$ comparison at $i-1$ and the $D$ comparison at $i$ imply
the $F$ comparison at $i$; then the $G$ comparison at $i-1$
implies the $G$ comparison at $i$. This finite induction proves
$G_k\ge\lambda G_{k-1}$ for EVERY $k\le N-3$.
Every D read lies on the actual increasing support through the
current rank; no post-mode identity with positive parts is used.
Flat or decreasing current slopes give C immediately.
This proves all-shift actual C from $W+1$ under (1).

The earlier reference tests and shorter-count/padding tests finish
the stated complete-C implications exactly as before.

\paragraph{Compatibility with the quadratic inverse reduction.}
At this new entrance $u=h+1$, the CI degree is $L=u+h=2u-1$.
Thus the exact derivative head through $2u$ covers EVERY triple
in (1). The stronger full-array inverse-kernel LC conjecture already
tested at degree at most 34 would directly supply this entrance if
proved for all CI tails; it is still OPEN. The proof above reduces
the sufficient entrance without assuming that conjecture.

\paragraph{An automatic actual band independent of entrance LC.}
For every $u,\alpha$, exact reference transport and the first-defect
payment already give ALL actual C shifts throughout
\[
                         W+1\le N\le u+r+2.
\]
The band is nonempty at $u\ge h-1$; neither (1) nor any density
LC assumption is needed for this band alone.

\paragraph{Scope.}
The general $u=h+1$ entrance, earlier arbitrary reference C,
general actual C/A2 and the ORIGINAL one-descent conjecture remain
OPEN. These are new universal implications and an unconditional
reference theorem, not a solution of any of those remaining claims.
